% !PRINTED 644
\weberpage{644}
But if $\mathfrak{p}=\mathfrak{p}'$, then $\mathfrak{P}=\mathfrak{p}$, and the formation of the norm yields $2f_1=f$.

Which of these two cases occurs therefore depends on whether the prime ideal $\mathfrak{p}$ admits the substitution $(r,r^{-1})$ or not, i.e. on whether $(r,r^{-1})$ is or is not contained in the group to which the prime ideal belongs. According to \S.\,172 this distinction is conditioned by whether there is any exponent $h$ for which the congruence
\[
p^h \equiv -1 \ (\mathrm{mod}\ m)
\]
is fulfilled, or whether there is no such exponent.

We summarise the result as follows:

2. \gesperrt{The natural primes $p$ different from $q$ split into two kinds $p_1, p_2$.}

\gesperrt{The first kind $p_1$ is characterised by the fact that for some exponent $h$}
\[
p_1^h \equiv -1 \ (\mathrm{mod}\ m),
\]
\gesperrt{and these primes, if they belong to the exponent $f$ for the modulus $m$ and $\mu=ef$ is set, split in the field $H_m$ into $e$ prime factors of degree $\frac12 f$.}

\gesperrt{Their prime factors are not further decomposable in $\Omega_m$ either.}

\gesperrt{The primes of the second kind $p_2$ are determined by the fact that none of their powers becomes congruent to $-1$ with respect to the modulus $m$, and they split in the field $H_m$ into $\frac12 e$ prime factors of degree $f$. Their prime factors are still decomposable in $\Omega_m$ into two factors each.}

For the primes of the first kind $f$ must certainly be an even number. If $m$ is odd, it also suffices, in order that $p$ be a prime of the first kind, that it belong to an even exponent; for then
\[
(p^{\frac12 f}-1)(p^{\frac12 f}+1)\equiv 0\ (\mathrm{mod}\ m).
\]
But the first of these factors $p^{\frac12 f}-1$ is not divisible by $m$, because otherwise $p$ would belong to the exponent $\frac12 f$. Consequently $p^{\frac12 f}+1$ must be divisible by $q$. But here the two factors $p^{\frac12 f}-1$ and $p^{\frac12 f}+1$ cannot both be divisible by $q$, because otherwise their difference, which is $=2$, would also be divisible by $q$; and consequently $p^{\frac12 f}+1$ is divisible by $m$.
